\documentclass{article}
\usepackage{fontspec} 
\usepackage{unicode-math}  % 先加载 unicode-math
\setmathfont{XITS Math}    % 再设置数学字体
%\usepackage[UTF8]{ctex} % 如果需要支持中文

\usepackage{anyfontsize}
\usepackage{xeCJK} % XeLaTeX 推荐 xeCJK，LuaLaTeX 也可用   
\usepackage{setspace}  % 添加这个包
\setstretch{1.2}      % 设置行距为1.5倍

\setCJKmainfont[
    Path = C:/USERS/11252/APPDATA/LOCAL/MICROSOFT/WINDOWS/FONTS/,  % 改用 Windows 字体目录
    UprightFont = {SourceHanSansCN-Light1.otf},  % 使用可变字体
    BoldFont = {SourceHanSansCN-Medium1.otf},
    LetterSpace = 2.0,  % 增加字间距
    WordSpace = 1.2,    % 增加词间距
    %BoldFeatures = {RawFeature={+wght=900}},
    %Scale=1
]{SourceHanSansCN}


\setmainfont[
    Path = C:/USERS/11252/APPDATA/LOCAL/MICROSOFT/WINDOWS/FONTS/,  % 改用 Windows 字体目录
    UprightFont = {SourceHanSansCN-Light1.otf},  % 使用可变字体
    BoldFont = {SourceHanSansCN-Medium1.otf},
    LetterSpace = 2.0,  % 增加字间距
    WordSpace = 1.2,    % 增加词间距
    %BoldFeatures = {RawFeature={+wght=900}},
    %Scale=1
]{SourceHanSansCN}

%\setCJKmainfont[
 %   Path = C:/USERS/11252/APPDATA/LOCAL/MICROSOFT/WINDOWS/FONTS/,
 %   UprightFont = {SourceHanSerifCN-Light-5.otf},
 %   BoldFont = {SourceHanSerifCN-Medium-6.otf},
%    Scale=1
%]{SourceHanSerifCN}

%\setmainfont{SourceHanSerif}  % 设置英文字体

% 处理冲突的命令
\let\oldbackepsilon\backepsilon
\let\backepsilon\relax
\let\oldeth\eth
\let\eth\relax
\let\olddigamma\digamma
\let\digamma\relax
\usepackage{float}

\usepackage[a4paper, left=0.1in, right=0.1in, top=0.05in, bottom=0.05in]{geometry} % 设置四边页边距为0.1英寸{geometry} % 设置页边距为0.5英寸
\usepackage{enumitem} % 用于自定义列表
\usepackage{amsmath}  % 数学公式支持
\usepackage{tikz}
\usetikzlibrary{arrows.meta, positioning}
\setlength{\parindent}{0pt} % 不缩进
\usepackage{etoolbox} % 用于列表操作
%调整类代码
\usepackage{comment}
\usepackage{tocloft} % 用于定制目录 样式
\usepackage{hyperref} %不需要目录盒子注释这
\usepackage{ifthen}
\usepackage{tikz}
\usetikzlibrary{arrows.meta}
\usepackage{titletoc}
\usepackage{graphicx}
\usepackage{float}

\usepackage{amssymb}  % 提供数学符号，包括\therefore
%目录设定
%快捷键 CMD + / 注释
%  \renewcommand{\cftdot}{} % 隐藏目录中的页码
%  \cftpagenumbersoff{section}
%  \cftpagenumbersoff{subsection}
%  \makeatletter
%  \renewcommand{\@pnumwidth}{0em}  % 设置页码宽度为0
%  \renewcommand{\@tocrmarg}{0em}   % 设置右边距为0
%  \makeatother
 

\usepackage{environ}

\newif\ifshowcontent  % 定义一个条件判断变量
\showcontenttrue    % 设置为 false，隐藏所有 question 环境的内容

\newcounter{question} % 题目计数器
\setcounter{question}{0}
\NewEnviron{question}{%
    \ifshowcontent
        \par\stepcounter{question}\textbf{\thequestion.} \BODY\par % 如果显示内容，则输出
    \fi
}

\NewEnviron{choices}{%
    \ifshowcontent
        \begin{enumerate}[label=\Alph*., leftmargin=*]
            \BODY
        \end{enumerate}\par
    \fi
}

\NewEnviron{correctanswer}{%
    \ifshowcontent
        \textbf{答案:}\BODY\par
    \fi
}



\newenvironment{explanation}
{\textbf{解析:}\par}
{\par}



% 新增考察方面命令
\newcommand{\checklist}[1]{
    \textbf{考察:} #1 \par % 在题目下方显示考察方面
    \addcontentsline{toc}{subsection}{题号\thequestion 考察: #1} % 将考察内容添加到目录
    \vspace{2em} % 添加1em的垂直空间
}

\graphicspath{{images/}} %统一


\begin{document}
 %\fontsize{22}{31}\selectfont
 \tableofcontents % 添加目录
\newpage
%\pagestyle{empty}




\section{考点1:}
\setcounter{question}{0}


\begin{question}
    Data analysts at an investment bank are using machine learning to detect money laundering activities. The method they use explores data and automatically finds patterns without requiring labels to train the algorithm. Which machine learning method are the analysts using?
\end{question}

\begin{choices}
    \item Classification.
    \item Clustering.
    \item Principal component analysis.
    \item Support vector machines.
\end{choices}

\begin{correctanswer}
    B
\end{correctanswer}

\begin{explanation}
    \textbf{A选项错误。}
    \begin{itemize}
        \item 分类方法需要标签来训练算法，属于监督学习
    \end{itemize}
    
    \textbf{B选项正确。}
    \begin{itemize}
        \item 聚类方法是无监督学习，可以自动发现数据模式而无需标签
    \end{itemize}
    
    \textbf{C选项错误。}
    \begin{itemize}
        \item 主成分分析是降维技术，不是专门用于模式发现的机器学习方法
    \end{itemize}
    
    \textbf{D选项错误。}
    \begin{itemize}
        \item 支持向量机是监督学习方法，需要标签进行训练
    \end{itemize}
\end{explanation}

\checklist{无监督学习方法}



\begin{question}
    A researcher at a financial supervisory agency is studying the rapid adoption of artificial intelligence (AI) and machine learning by financial institutions. The researcher prepares a report on potential benefits and risks to financial markets from these technologies. Which conclusion would the researcher correctly reach?
\end{question}

\begin{choices}
    \item AI and machine learning help improve data quality by automatically highlighting potential errors to flag them to users.
    \item AI and machine learning minimize the risk of creating "black boxes" in decision making and potential system-wide stress.
    \item While AI and machine learning lower trading costs, they will not improve financial securities' price discovery process.
    \item While AI and machine learning may deliver customized services to investors, they increase customers' fees and borrowing costs.
\end{choices}

\begin{correctanswer}
    A
\end{correctanswer}

\begin{explanation}
    \textbf{A选项正确。}
    \begin{itemize}
        \item AI和机器学习可以通过自动标记潜在错误来帮助提高数据质量
    \end{itemize}
    
    \textbf{B选项错误。}
    \begin{itemize}
        \item AI和机器学习实际上可能增加决策中的"黑盒"风险和系统性压力
    \end{itemize}
    
    \textbf{C选项错误。}
    \begin{itemize}
        \item AI和机器学习不仅降低交易成本，还能改善金融证券的价格发现过程
    \end{itemize}
    
    \textbf{D选项错误。}
    \begin{itemize}
        \item AI和机器学习通常降低客户费用和借贷成本，而不是增加
    \end{itemize}
\end{explanation}

\checklist{AI金融应用影响}



\begin{question}
    A senior risk manager at a global bank is proposing greater use of artificial intelligence (AI) and machine learning (ML) systems in the bank's risk management practices. In a presentation to the CEO and CFO, the manager provides examples of how these approaches can help mitigate different types of risk exposures. Which example most appropriately describes how AI or ML can mitigate the given risk?
\end{question}

\begin{choices}
    \item ML can help mitigate operational risk by replacing human judgment-based processes with complex algorithms for lending decisions.
    \item ML can help mitigate model risk by detecting changes in trading patterns that diverge from modeled estimates.
    \item AI can help mitigate credit risk by analyzing customer account documentation for instances of missing or fraudulent entries.
    \item AI can help mitigate market risk by identifying a large position in one asset that traders can take to replicate the performance of a portfolio of smaller, more liquid positions.
\end{choices}

\begin{correctanswer}
    C
\end{correctanswer}

\begin{explanation}
    \textbf{A选项错误。}
    \begin{itemize}
        \item 用复杂算法替代人工判断可能增加操作风险，而不是缓解
    \end{itemize}
    
    \textbf{B选项错误。}
    \begin{itemize}
        \item 检测交易模式变化是市场风险缓解，不是模型风险缓解
    \end{itemize}
    
    \textbf{C选项正确。}
    \begin{itemize}
        \item AI分析客户账户文件以发现缺失或欺诈条目是缓解信用风险的恰当应用
    \end{itemize}
    
    \textbf{D选项错误。}
    \begin{itemize}
        \item 识别大仓位来复制小仓位组合表现不是市场风险缓解的有效方法
    \end{itemize}
\end{explanation}

\checklist{AI风险缓解应用}


\begin{question}
    A key difference between low-inflation and high-inflation regimes is that:
\end{question}

\begin{choices}
    \item High-inflation regimes are increasingly sensitive to relative price shocks.
    \item High-inflation regimes generally tend to be self-equilibrating (self-stabilizing).
    \item Low-inflation regimes tend to see a constant level of spillovers across individual price changes.
    \item In a low-inflation regime, lower inflation volatility is due to a decline in the volatility of individual price changes.
\end{choices}

\begin{correctanswer}
    A
\end{correctanswer}

\begin{explanation}
    \textbf{A选项正确。}
    \begin{itemize}
        \item 高通胀制度对相对价格冲击越来越敏感，这是与低通胀制度的关键区别
    \end{itemize}
    
    \textbf{B选项错误。}
    \begin{itemize}
        \item 高通胀制度通常不是自我平衡的，而是不稳定的
    \end{itemize}
    
    \textbf{C选项错误。}
    \begin{itemize}
        \item 低通胀制度中个体价格变化的溢出效应不是恒定的
    \end{itemize}
    
    \textbf{D选项错误。}
    \begin{itemize}
        \item 低通胀制度中通胀波动性降低不是由于个体价格变化波动性下降
    \end{itemize}
\end{explanation}

\checklist{通胀制度差异}





\newpage


\end{document}